\scanpage{35}
\mspage{9}
définit donc (avec la paire $\{z,z'\}$) une métrique sphérique, ce qui sera OK. On fera attention\inserted{a priori}, que $\sigma$ n'est pas unique, il y a une ambiguïté d'ordre $2$, on peut le remplacer par $\sigma'=\varepsilon.\sigma$, où $\varepsilon$ est l'élément d'ordre $2$ de $T$. On trouve en fait que $\{\sigma,\sigma'\}$ est formé d'une inversion et d'un antipodisme.

\uline{Dém.} OPS $X=$ sphère de Riemann, $\{z,z'\}=\{\infty,0\}$, alors
\[
T=\{z\longmapsto az\mid a\in\mathbb U\text{ i.e. }a\in\mathbb C,\ |a|=1\text{ i.e. }a\bar a=1\}
\]
les \deleted{\unsure{anti}} autom. conformes anti-hol. de $\mathbb C^*$ sont
\diagram{\begin{array}{lll}
z\longmapsto a\bar z &(a\in\mathbb C^*)&\text{ceux qui fixent }0,\infty\\
z\longmapsto b/\bar z &(b\in\mathbb C^*)&\longrightarrow\ \text{intervertissent }0,\infty
\end{array}}
On voit tout de suite que ce sont les \uline{\unsure{seconds|seuls}} qui commutent à $T$. \deleted{Une telle} Si
\[
\sigma(z)=b/\bar z,\quad\text{on a}\quad\sigma^2(z)=b/\bar b\ z,
\]
donc $\sigma$ involutif ssi $b/\bar b=1$ i.e. $b\in\mathbb R^*$. Les pts fixes de $\sigma$ sont donnés par
\[
z\bar z=b
\]
\deleted{\unsure{les}} donc $\sigma$ est une inversion resp. un antipodisme, suivant que \deleted{$b\in\mathbb R$} $b>0$ ou $b<0$. Si $\sigma(z)=b/\bar z$, on a $\sigma'(z)=-b/\bar z$, \deleted{\unsure{donc}} donc si \deleted{\unsure{l'inv.}}\inserted{$\sigma$} est une inversion, $\sigma'$ est un antipodisme, et vice versa. OPS les notations choisies de telle façon que $b>0$, i.e. $\sigma$ inversion, donc $\sigma$ définit une métrique sphérique ayant \unsure{le|leur} grand cercle défini par
\[
\unsure{C|E}=\{z\in\mathbb C\mid z\bar z=b\}.
\]
% archive foliation: 37
